\chapter*{How to Read This Book: Four Paths}
\addcontentsline{toc}{chapter}{How to Read This Book: Four Paths}
\markboth{How to Read This Book}{How to Read This Book}

Every chapter has the same shape, so the reader knows what to expect and may skip ahead
without losing the way.

\begin{center}
\small
\begin{tabular}{@{}p{3.1cm}p{4.6cm}p{4.4cm}@{}}
\toprule
\textbf{Path} & \textbf{What is read} & \textbf{What you get} \\
\midrule
Younger reader (10-14) & the \emph{For younger readers} box, the Summary, one worked
  example & the core idea without notation; you know what is being computed \\
Undergraduate (S1) & the whole body of the chapter, the Computational lab, the Exercises;
  skip the researcher box & definitions, proofs, the ability to compute and to check \\
Instructor & the \emph{Note for instructors} box, the source map, the Exercises & a plan
  for the meeting, misconceptions to anticipate, blackboard material \\
Researcher & the \emph{For researchers} box, the source map, Appendices~A and~B & the
  limits of the claims, primary sources, open questions, consistent notation \\
\bottomrule
\end{tabular}
\end{center}

\medskip
The \textsc{level} mark at the start of every section shows how hard the reading is:
\emph{basic} can be read without any earlier chapter other than the ones named;
\emph{intermediate} demands the algebra and the notation of this book;
\emph{advanced} may be skipped on a first reading without breaking the flow of the
chapter.

\medskip
\noindent\textbf{Six ways of learning, six doors in.} People do not learn in the same
way. So every chapter carries more than one route to the same idea, and no route is
compulsory: pick the one that makes you understand fastest, then go back to the body of
the chapter to see why that way works.

\begin{center}
\small
\begin{tabular}{@{}p{3.4cm}p{4.3cm}p{4.6cm}@{}}
\toprule
\textbf{Type} & \textbf{Look for this in every chapter} & \textbf{How to use it} \\
\midrule
Visual & 2-3 numbered figures (in the List of Figures), tables, and the box colors & read
  the figure before its text; the figure carries the same numbers as the computation in
  the text \\[2pt]
Auditory & the \emph{Read it aloud} box (indigo) & read the text aloud or record your
  voice; repeat the passages that have a rhythm \\[2pt]
Kinesthetic & the \emph{Do it by hand} box (green) & do it with real objects: paper,
  cards, the floor, money, or your own body \\[2pt]
Concept-first & definitions, theorems, proofs, the \emph{For researchers} box & do not
  use a formula before you know where it comes from; the proof is where that origin is
  written \\[2pt]
Procedural memorizer & the \emph{Step-by-step recipe} box (brown) and the Exercises &
  memorize the order of the recipe cards, then apply it to different numbers until it is
  automatic \\[2pt]
Problem solver & the \emph{Challenge} box (purple) and the starred exercise items & take
  the challenge first; the hints are tiered, do not open them all at once \\
\bottomrule
\end{tabular}
\end{center}

The four paths below set \emph{how much} is read; the six types above set \emph{which
door} you come in through. The two may be combined: a visual reader who pursues the
researcher path still reads every proof, only starting from the figure.

Three learning paths that have already been tried, from the lightest:

\begin{enumerate}
  \item \textbf{Curiosity path (8-12 weeks).} Read the orange box, the Summary, and one
        worked example from every chapter, in order from Chapter~1 to Chapter~32. No
        exercises. The aim is one thing: to know the map of the terrain.
  \item \textbf{Student path (2 semesters).} First semester: Chapters~1-10 and
        Chapter~17. Second semester: Chapters~11-16, then 18-32. That is twenty-one
        chapters, more than an ordinary semester; the way to trim it without breaking the
        prerequisites is to postpone Chapters~31-32 (cryptography and privacy) to an
        elective course, so that nineteen chapters remain. Do the Computational lab of
        every chapter; at least four exercise items per chapter. This order keeps the
        promise about prerequisites: Chapter~18 demands Chapters~13 and~17, so the three
        must not be far apart; Chapter~21 demands Chapters~10 and~20; Chapter~23 demands
        Chapter~12.
  \item \textbf{Machine-learning security researcher path (9 weeks).} Nineteen chapters,
        not seven: Chapters~3, 4, 9, 10, 12, 13, 18, 19, 20, 22, 23, 25, then 26-30,
        closing with Chapters~31 and~32. Chapters~13, 18, 22, and~23 are included
        because each is a prerequisite of Chapter~20 (hypothesis testing), Chapter~29
        (adversarial), and Chapter~25 (neural networks); Chapters~31 and~32 complete the
        cryptography and privacy side that the title of the book demands. Reading this
        path in six weeks means skipping the prerequisites and losing the reason why the
        claims hold.
\end{enumerate}

Symbols are used consistently throughout the book: $\mathbb{N}$, $\mathbb{Z}$,
$\mathbb{Q}$, $\mathbb{R}$ for the number sets; $\mathbf{x}$ for a vector and $x_i$ for
its components; $\|\mathbf{x}\|$ for the Euclidean norm; $\nabla f$ for the gradient;
$\varepsilon$ and $\delta$ always small positive numbers; $\mathbb{E}[X]$ for the
expected value; $\mathcal{L}(V,W)$ for the space of linear transformations; $\top$ for
the transpose. Appendix~A contains the complete list.
